%% 3.1.9 Resumo em língua portuguesa
%% Elemento obrigatório (Figura 14).
%% Consiste na apresentação sucinta dos pontos relevantes do texto,
%% em um único parágrafo. O resumo deve conter entre 150 e 500 pala-
%% vras e fornecer uma visão rápida e clara dos objetivos, da metodologia,
%% dos resultados e das conclusões do trabalho. Na elaboração do resumo,
%% deve-se usar o verbo na voz ativa, na terceira pessoa do singular.

%% Fonte -> TNR ou Arial, corpo 12.
%% A palavra RESUMO deve aparecer em letras maiúsculas
%% e em negrito.
%% O uso de itálico é permitido em palavras estrangeiras.
%% O uso de letras maiúsculas nas palavras-chave
%% restringe-se ao início da palavra, em nomes próprios
%% e siglas, se for o caso.

%% Alinhamento -> A palavra RESUMO deve estar localizada na margem
%% superior da folha e centralizada, e a referência, alinhada
%% à margem esquerda;
%% O alinhamento é justificado para o texto do resumo,
%% que inicia com parágrafo, e para as palavras-chave.

%% Espaçamento -> A palavra RESUMO deve ser separada da referência por
%% duas linhas em branco de 1,5;
%% Espaço 1 na referência e no resumo e, nas palavras-
%% chave, espaço 1,5.

%% Formato do papel,
%% orientação e margens -> Conforme especificado na seção 1.1.

%% Pontuação -> As palavras-chave devem ser separadas por ponto e
%% terminadas por ponto.

\begin{resumo}
\begin{SingleSpace}

\noindent
\begin{flushleft}
\entradaAutor{}. \textit{\imprimirtitulo}. \the\year. \pageref{LastPage} f. Trabalho de Conclusão de Curso (Graduação em Engenharia de Computação) - Instituto Politécnico, Universidade do Estado do Rio de Janeiro, Nova Friburgo, \the\year.
\end{flushleft}
\vspace{\onelineskip}

\setlength{\parindent}{1.3cm}

Este trabalho apresenta o desenvolvimento do Onmarmoraria, um sistema ERP voltado para marmorarias, com o objetivo de centralizar a gestão de clientes, produtos e orçamentos em um painel web. O sistema é \textit{multi-tenant}: uma única instalação atende várias marmorarias contratantes, cada uma com conta, usuários e dados próprios e isolados das demais. A metodologia envolveu levantamento de requisitos do setor, modelagem de dados, implementação de uma API REST em Go com autenticação JWT, \textit{frontend} em Next.js e testes automatizados e manuais. A arquitetura foi desenhada para execução em nuvem, com suporte a implantação em máquina virtual (IaaS) ou VPS e a orquestração via Kubernetes, cujos manifestos já integram o repositório. A validação deste trabalho, porém, foi conduzida em ambiente Docker local, reproduzindo os mesmos contêineres previstos para produção, dado o caráter de MVP acadêmico e a janela de custo do projeto. O armazenamento de mídias foi implementado com um \textit{backend} de disco local, mantendo o mesmo contrato de URLs assinadas usado por serviços de nuvem (\textit{Google Cloud Storage} ou um serviço compatível com o Amazon S3, como o Cloudflare R2), de modo que a troca de \textit{backend} não exija alterações de código; as imagens Docker da API são publicadas no GHCR por meio de um \textit{pipeline} de CI/CD no GitHub Actions.

\vspace{\onelineskip}
\noindent Palavras-chave: ERP; Marmoraria; Orçamentos; Computação em nuvem; Conteinerização.

\end{SingleSpace}
\end{resumo}